\section{Conclusion and Discussion}

In this paper, we presented AutoShotV2, a compact refinement framework for short-video shot boundary detection. By freezing the AutoShot backbone and optimizing a two-branch classification head with Focal Loss and many-hot supervision, the method addresses the extreme class imbalance characteristic of short videos. Temperature scaling and Gaussian smoothing provide the calibrated decision pipeline used at deployment. AutoShotV2 raises SHOT F1 from \PaperAutoShotShotFOne{} for the reproduced AutoShot baseline to \PaperASVTwoShotFOne, the best result among the evaluated methods, without repeating the computationally expensive neural architecture search.

Despite these advancements, several challenges remain. Cross-domain analysis revealed a noticeable performance degradation on diverse open-world datasets like ClipShots, primarily due to missed gradual transitions. This emphasizes that SBD systems should treat calibration, loss balancing, and temporal windows as domain-sensitive parameters rather than fixed constants. Fast camera motion and game-like visual effects can still trigger false positives, indicating a limit to purely visual features. In the future, we aim to explore domain-adaptive calibration techniques and multimodal SBD architectures that integrate audio cues to further improve cross-domain robustness.
